\section{Varying the coarse field while holding the fine field fixed}
\label{sec:compensated-variation}
Differentiating a packet directly in its coarse field would also
change the fine configuration represented by its fluctuation variable.
The resulting derivatives would meet the fine masks and branch choices.
We instead make the actual compensated change of variables. It is a
global bijection on the fluctuation space near the chosen coarse
configuration, and agrees with the required logarithmic adjustment
where the protected integrand can be nonzero.

\subsection{A globally defined adjustment of one logarithm}
We first specify the logarithm on nearby ambient matrices. As in
Section~\ref{sec:predictor-branch}, regard $\mathcal M=M_N(\mathbb C)$
as a real inner-product space with
$\langle A,B\rangle=\operatorname{Re}\operatorname{tr}(A^*B)$, and split
it as the colour space $\mathfrak g$ and its orthogonal complement
$\mathfrak g^\perp$. The map $(u,n)\mapsto\exp(u)\exp(n)$ on
$\mathfrak g\times\mathfrak g^\perp$ has derivative $(u,n)\mapsto u+n$
at $(0,0)$, which is invertible. By the inverse function theorem there
is a neighbourhood $O$ of $(0,0)$ in $\mathfrak g\times\mathfrak g^\perp$
such that this map is a diffeomorphism of $O$ onto an open neighbourhood
$V$ of the identity $1$ in $\mathcal M$. Its inverse
$V\to O$ is smooth; let $\operatorname{Log}_*\colon V\to\mathfrak g$
be the colour component of this inverse.
In particular, $\operatorname{Log}_*$ is colour-valued even when its
matrix argument is not a group element.

For a group element $g$ near the identity, $\operatorname{Log}_*(g)$
is the local group logarithm. This uses one fact from the theory of
matrix Lie groups, already used in Section~\ref{sec:predictor-branch}:
by Cartan's closed-subgroup theorem in the form of B.~C. Hall,
\emph{Lie Groups, Lie Algebras, and Representations}, 2nd ed.,
Springer, 2015, Theorem~3.42, there is $\theta>0$ such that every
$g\in G$ with $\|g-1\|<\theta$ equals $\exp X$ for some
$X\in\mathfrak g$, with $X\to0$ as $g\to1$. For such $g$ close
enough to $1$, the pair $(X,0)$ lies in $O$ and is mapped to $g$, so
$g\in V$ and, the map being injective on $O$, $\operatorname{Log}_*(g)=X$.
Without this theorem one would know only that $g$ is a matrix near
$1$, and its local preimage could have a nonzero
$\mathfrak g^\perp$ component.

Choose fixed $d_*,A>0$ so small that
\[
 \operatorname{Log}_*(\exp u)=u\quad(u\in\mathfrak g,\ \|u\|\le2d_*)
\]
and this extension is smooth at $\exp(u)M$ whenever
$\|u\|\le2d_*$ and $\|M-1\|\le A$. Also require
\[
 \exp\operatorname{Log}_*(\exp(u)g)=\exp(u)g
 \qquad(u\in\mathfrak g,\ \|u\|\le d_*,\ g\in G,\ \|g-1\|\le A).
\]
This identity is asserted only for group elements $g$: for a matrix
$X$ outside $G$ the left-hand side lies in $G$ and the right-hand side
does not. Such constants exist by continuity of multiplication
and the exponential at $(0,1)$, after choosing a smaller product
neighbourhood inside the given logarithm domain; for the last identity
note that $\exp(u)g\in G$ is close to $1$, so the previous paragraph
applies to it. The constants are chosen before the coupling threshold
and before the particular cutoff. The fixed chart radius $d_*$ is not
the growing scale $d=1+B+(\log h)^2$ of \eqref{eq:warning-scales}.

Using the same smooth cutoff as before, define
\[
 \psi(M,u)=\chi(\|u\|^2/d_*^2)\{\operatorname{Log}_*(\exp(u)M)-u\},
 \qquad F_M(u)=u+\psi(M,u).
\]
Here group elements are represented by their unitary matrices in the
logarithm expression. The correction is zero for $\|u\|\ge2d_*$.
On the closed ball it is smooth in $(M,u)$; outside the ball it
vanishes locally, so these formulas define a smooth map everywhere
for $M$ sufficiently near $1$. At $M=1$ the correction and all its
$u$ derivatives vanish.

Compactness of $\|u\|\le2d_*$ now supplies one neighbourhood of
$M=1$ on which, for every $u$,
\begin{equation}\label{eq:compensator-small}
 \|\psi(M,u)\|<d_*/2,\qquad \|D_u\psi(M,u)\|<1/2.
\end{equation}
Indeed these inequalities hold at $(1,u)$, and their continuous
left-hand sides are uniformly small on that compact set; off it they
are zero. Integration of the derivative on line segments shows that
$\psi(M,\cdot)$ is $1/2$-Lipschitz. Thus $F_M$ is injective, since
$F_M(u)=F_M(v)$ implies $\|u-v\|\le\|u-v\|/2$.
For every $z$, iteration of $u\mapsto z-\psi(M,u)$ is a contraction
on the complete colour space. Its iterates are Cauchy because successive
distances form a geometric series; their limit is its unique fixed
point and solves $F_M(u)=z$. Hence $F_M$ is surjective as well.

Its derivative is $1+D_u\psi$, which is invertible by the convergent
geometric series. The determinant
\[
 p(M,u)=\det(1+D_u\psi(M,u))
\]
is positive for all $u$ after a further uniform restriction on $M$:
it equals one at $M=1$, is continuous on the compact $u$ ball, and
is one outside it. Thus $F_M$ is a smooth bijection with positive
Jacobian. Its local smooth inverses agree by injectivity, giving a
global smooth inverse.

For an actual group element $g$ near $1$ and $\|u\|\le d_*$, the
cutoff equals one. Consequently
\begin{equation}\label{eq:compensator-exact}
 \exp(F_g(u))=\exp(u)g.
\end{equation}
The map is global, but this exact group identity is used only within
that smaller ball.

\subsection{The coherent coarse and fluctuation variables}
Fix a sector with coarse field $V_0$, and keep its seeds, forest,
tags, frozen fine bonds and active sets fixed. In the affine source
replace the clipped coarse guide by its unmasked path-field formula;
let $\widetilde\ell(V)$ be the resulting linear source and put
\[
 \nu(V)=-P(V)^*\widetilde\ell(V)
\]
in physical colour coordinates. This specifies the extension used
in the change of variables. At a valid sector, every selected row
has a clean neighbourhood. The clipping is then inactive at every
read of that row, as proved in the guide comparison. Its unmasked
source equals the original source, and hence $\nu(V_0)=\mu$.
The same statement holds at any other valid sector. Away from valid
sectors the unmasked expression is retained as its definition.

Let $g_\beta=\beta^{-1/2}>0$ and identify a counting Gaussian
coordinate $z$ with its colour vector on each fine bond. The extended
fluctuation field is
\[
 u_0=\nu(V_0)+g_\beta z.
\]
For a nearby coarse field $V$, define separately at each fine bond
\begin{equation}\label{eq:compensator-course}
 M_i(V)=\overline V_{0,i}\overline V_i^{-1},\quad
 u_i(V)=F_{M_i(V)}(u_{0,i}),\quad
 z_i(V)=g_\beta^{-1}\{u_i(V)-\nu_i(V)\}.
\end{equation}
The replica is the prescribed one. It follows identically that
$\nu(V)+g_\beta z(V)=u(V)$, and at $V=V_0$ the map is the identity
on $z$.

The fine field on an active bond is $\exp(u_i)\overline V_i$;
on an inactive bond it is the fixed original fine field. Suppose
$8a\tau\le d_*/2$ and the varied active fluctuation satisfies
$\|u_i(V)\|\le8a\tau$. By \eqref{eq:compensator-small},
\[
 \|u_{0,i}\|\le\|u_i(V)\|+\|\psi(M_i(V),u_{0,i})\|<d_*.
\]
Equation~\eqref{eq:compensator-exact} therefore gives
\begin{equation}\label{eq:compensator-held-fine}
 \exp(u_i(V))\overline V_i
 =\exp(u_{0,i})\overline V_{0,i}.
\end{equation}
Inactive bonds agree by definition. On the protected support all fine
reads are consequently fixed, including their masks, likelihoods and
selected branch tags. The larger $8a\tau$ allowance is deliberate: at
a point where the broader cutoff germ is not identically zero,
continuity first gives $\|u_i(V_0)\|\le4a\tau$, then gives the strict
$8a\tau$ bound on a neighbourhood. A finite set of protected bonds has
one common such neighbourhood. No condition is imposed on the
unprotected Gaussian coordinates.

\subsection{Uniform derivative bounds for the compensation}
Fix the derivative order $k$. On
$\{1\}\times\{u:\|u\|\le2d_*\}$, all derivatives in $M$ of
$\psi$ and $p$ through order $k$ have fixed bounds, by smoothness and
compactness. Outside this set their $u$-sections are respectively
zero and one as functions of $M$. Hence these bounds hold for every
$u\in\mathfrak g$, with constants independent of the Gaussian
coordinate. Moreover $\psi(1,u)=0$ and $p(1,u)=1$.

Suppose the coarse variation $V(y)$ has the link derivative bounds
of Section~\ref{sec:normalization-jets}, with $V(y_0)=V_0$ and
scale $v(x)$ of upper exponent $-1/2$. Inversion, replication and
multiplication by the fixed unitary replica show that $M_i(V(y))$
has the same bound. In a chain-rule term of positive order $r\le k$,
the product of derivatives of $M_i$ has total scale $v^r$; its
amplitude is bounded by a constant depending only on $k$ and the
representation. There are finitely many partitions of $r$. Thus,
for every $u_{0,i}$,
\[
 \|D^r_y\psi(M_i(V(y_0)),u_{0,i})\|\le C_kv^r,
 \qquad
 \|D^r_y p(M_i(V(y_0)),u_{0,i})\|\le C'_kv^r
 \quad(1\le r\le k).
\]
Their zero-order values are zero and one. Since $v\to0$, enlarge
constants and obtain
\begin{equation}\label{eq:compensator-jets}
 \|D^r_y(u_i(y)-u_{0,i})|_{y=y_0}\|\le s(x),\qquad
 \|D^r_y p(M_i(V(y)),u_{0,i})|_{y=y_0}\|\le D(x)^r
 \quad(0\le r\le k),
\end{equation}
where $s\ge0$ has upper exponent $-1/2$, and $D\ge0$ has upper
exponent zero. For the first bound use $C_kv$ for positive orders
and zero at order zero; for the second take
$D=\max(1,C'_k)v$, increasing the fixed constant if necessary.
For $r\ge1$, $\max(1,C'_k)^r\ge C'_k$. At order zero the second
bound is equality, with the convention $D^0=1$.
These are precisely the compensated increment bounds used in
Section~\ref{sec:haar-chart}. They are uniform over all Gaussian
samples; they do not require a bound on $z_i(y)$ or on the derivatives
of the mean extension.

The incoming link-variation premise itself holds for the prescribed
scaled exponential variations. For example, let
$V_i(y)=\exp(\epsilon L_i y)V_{0,i}$, where the linear maps
$L_i$ have uniformly bounded norm and $y$ stays in a fixed bounded
set. The derivatives of the matrix exponential are uniformly bounded
on the resulting fixed colour ball. The linear chain rule contributes
$\epsilon^r$ at order $r$. Right multiplication by $V_{0,i}$ is
an isometry. Enlarging one fixed constant therefore gives the link
bound with $v=C_k\epsilon$. In particular the required scale holds
when $\epsilon$ has upper exponent $-1/2$, as for the prescribed
inverse-square-root coupling normalization. Finite collections of
marked directions are handled by taking the maximum of their norms.

\subsection{The exact finite-dimensional Jacobian}
For a fixed nearby $V$, the map in
\eqref{eq:compensator-course} acts separately on the colour coordinate
of every fine bond. On a finite set $P$ of bonds it is conjugate,
by the fixed coordinate regrouping, to the product of maps
\[
 z_i\longmapsto g_\beta^{-1}
       \{F_{M_i(V)}(\nu_i(V_0)+g_\beta z_i)-\nu_i(V)\}.
\]
Each is a bijection. Its derivative is exactly
$D F_{M_i(V)}(u_{0,i})$: the two scalar factors $g_\beta^{-1}$ and
$g_\beta$ cancel. The derivative of the product is block diagonal,
so its determinant is
\begin{equation}\label{eq:compensator-jacobian}
 \prod_{i\in P}p(M_i(V),u_{0,i})>0.
\end{equation}
A fixed regrouping of counting coordinates conjugates the derivative
and does not change its determinant. This remains true if the
regrouped spaces are temporarily given different finite-dimensional
norms. For the empty set the map is the identity of the zero-dimensional
space and the product is one.

The ordinary change-of-variables formula therefore gives, for every
nonnegative measurable test $f$ on these coordinates (and hence for
every integrable test),
\[
 \int f(z')\,dz'
 =\int f(z(V))\prod_{i\in P}p(M_i(V),u_{0,i})\,dz.
\]
Together with \eqref{eq:compensator-held-fine} and
\eqref{eq:compensator-jets}, this supplies the coherent protected
variation and its exact Jacobian. Differentiating a Gaussian density
under this substitution still requires its own covariance and mean
score bounds; those do not follow from this Jacobian calculation.
